Fix \(p\in\Delta_K\) and \(0<\gamma<1\).  Every \(B_{k\ell}>0\), hence \(q_k>0\), and \(\texttt{lem:exact-hit-factorization}\) gives, for \(z\in\mathcal S_k\),
\[
r_z:=\Pr(Z_c=z)=q_k/M_k>0. \tag{66}
\]
Set \(\widehat r_z=N_z/C\) and
\[
A_{C,z}=C^{-1}\sum_{c=1}^C\mathbf1\{Z_c=z\}\{W_c(z)-\mu_z\}.
\]
Boundedness gives uniformly in \(P\)
\[
\mathbb EA_{C,z}^2\leq C^{-1},\qquad
\mathbb E(\widehat r_z-r_z)^2\leq C^{-1}. \tag{67}
\]
Since there are finitely many positive-probability target cells, Chebyshev and a union bound imply
\[
\max_z|A_{C,z}|=O_P(C^{-1/2}),\quad
\max_z|\widehat r_z-r_z|=O_P(C^{-1/2}),\quad
\min_z\widehat r_z\geq\tfrac12\min_zr_z
\]
with probability tending to one, uniformly in \(P\).  On that event, expanding every cell ratio gives
\[
\sqrt C\{\widetilde U_C-U(P)\}
=C^{-1/2}\sum_{c=1}^C\xi_c(P)+o_P(1),\qquad
\xi_{c,k}=q_k^{-1}\mathbf1\{Z_c\in\mathcal S_k\}
\{W_c(Z_c)-\mu_{Z_c}\}. \tag{68}
\]
The remainder is \(o_P(1)\) after scaling because it is a finite sum of products of the two \(O_P(C^{-1/2})\) errors in (67).  This is the canonical-gradient expansion in \(\texttt{thm:efficient-exact-hit-bernoulli}\).

Conditioning on cells gives
\[
\mathbb E\xi_{c,k}=0,\qquad
\mathbb E\xi_{c,k}^2=\tau_k^2(P)/q_k,\qquad
\mathbb E(\xi_{c,j}\xi_{c,k})=0\quad(j\ne k). \tag{69}
\]
Thus the covariance is \(\Sigma_B^{(K)}(P)\).  To prove uniformity, take any sequence \((P_C)\) with \(\min_k\tau_k^2(P_C)\geq\underline\tau^2\).  The schedule space \([0,1]^{n\times\mathcal Z_n}\) is compact metric, so \(\texttt{lem:weak-compactness-compact-laws}\) supplies from every subsequence a further subsequence with \(P_C\Rightarrow P_\infty\).  The coordinate maps \(W(z)\) and \(W(z)^2\) are bounded continuous, whence
\[
\Sigma_B^{(K)}(P_C)\longrightarrow\Sigma_B^{(K)}(P_\infty)=:\Sigma_\infty. \tag{70}
\]
The limit is positive definite by the variance lower bound.  For fixed \(a\), \(a'\xi_c(P_C)\) is uniformly bounded, so
\[
\mathbb E_{P_C}e^{\mathrm ia'\xi_c(P_C)/\sqrt C}
=1-\frac{a'\Sigma_B^{(K)}(P_C)a}{2C}+O(C^{-3/2}), \tag{71}
\]
uniformly along the sequence.  Independence, (68), and (70)--(71) give
\[
\sqrt C\{\widetilde U_C-U(P_C)\}\Rightarrow N(0,\Sigma_\infty). \tag{72}
\]

For studentization, on the regular branch put
\[
\widehat\xi_{c,k}=\widehat q_k^{-1}\mathbf1\{Z_c\in\mathcal S_k\}
\{W_c(Z_c)-\overline W_{Z_c}\}.
\]
The within-cell identity
\[
\sum_{c:Z_c=z}(W_c(z)-\overline W_z)^2
=\sum_{c:Z_c=z}(W_c(z)-\mu_z)^2-N_z(\overline W_z-\mu_z)^2
\]
and (67) show uniformly that centering by \(\overline W_z\) instead of \(\mu_z\) changes empirical second moments by \(o_P(1)\).  Bounded laws of large numbers and \(\widehat q_k\to q_k\) then give
\[
\widehat\Sigma_C^{(K)}-\Sigma_B^{(K)}(P)\to_P0 \tag{73}
\]
uniformly over the displayed class.

Let \(e_{ij}=e_i-e_j\), define
\[
T_C=\max_{i<j}\frac{\sqrt C|e_{ij}'(\widetilde U_C-U)|}
{(e_{ij}'\widehat\Sigma_C^{(K)}e_{ij})^{1/2}},
\]
and let \(c_C(\gamma)\) be the conditional \((1-\gamma)\)-quantile of the analogous Gaussian maximum with covariance \(\widehat\Sigma_C^{(K)}\).  Every limiting contrast variance is at least \(\underline\tau^2(q_i^{-1}+q_j^{-1})>0\).  Hence (72)--(73) imply
\[
T_C\Rightarrow\max_{i<j}\frac{|e_{ij}'G_\infty|}
{(e_{ij}'\Sigma_\infty e_{ij})^{1/2}},\qquad G_\infty\sim N(0,\Sigma_\infty). \tag{74}
\]
This Gaussian maximum has a continuous distribution: its contrast vector has a positive density on the \((K-1)\)-dimensional contrast space and max-norm spheres have Lebesgue measure zero.  Thus plug-in Gaussian quantiles are continuous and simultaneous pairwise coverage tends to \(1-\gamma\).  If uniform coverage failed, a violating sequence and the weakly convergent subsequence above would contradict (74).  This proves the displayed uniform liminf.

On simultaneous contrast coverage, a policy certified above \(k\) is truly above it, and a policy certified below \(k\) is truly below it.  Therefore the inverted lower rank endpoint is at most \(r_k^-(U)\) and its upper endpoint is at least \(r_k^+(U)\), simultaneously for all \(k\).  This deterministic implication uses no separation and covers all partial ties.  Finally
\[
\Pr(N_z=0)=(1-q_k/M_k)^C\leq e^{-Cq_k/M_k}.
\]
A union bound gives the stated fallback probability.  The fallback returns all ranks, so it cannot reduce coverage.